\documentclass[12pt,a4paper]{ctexart}

\usepackage{geometry}
\geometry{top=2.5cm,bottom=2.5cm,left=2.5cm,right=2.5cm}
\usepackage{amsmath,amssymb,mathtools}
\usepackage{enumitem,array,booktabs,pifont}
\usepackage{hyperref,fancyhdr}
\usepackage{bookmark}

\pagestyle{fancy}
\fancyhf{}
\fancyhead[C]{\textbf{数学 · 限时模拟卷一}}
\fancyfoot[C]{\thepage}

\title{\textbf{数学限时模拟卷一}\\[5pt]\large 满分150分\quad 时间120分钟}
\date{}

\begin{document}
\maketitle

%=================== 各档次策略 ===================
\section*{各档次答题策略}

\begin{center}
\begin{tabular}{p{2.5cm}|p{4cm}|p{4cm}|p{4cm}}
\textbf{目标} & \textbf{必做题（保分）} & \textbf{尽力题（冲刺）} & \textbf{建议放弃} \\ \hline
\textbf{90分} & 单选1--6、多选9--10、填空13--14、解答17--18 & 解答19第一问、20第一问 & 单选7--8、多选11--12、填空15--16、解答21--22后两问 \\ \hline
\textbf{120分} & 单选1--7、多选9--11、填空13--15、解答17--20 & 单选8、填空16、解答21、22第一问 & 多选12完全放弃、解答22第二问 \\ \hline
\textbf{140+分} & 全部题目 & 全做 & 无（每分必争） \\
\end{tabular}
\end{center}

\subsection*{时间分配建议}
\begin{center}
\begin{tabular}{c|c|c|c}
\textbf{题型} & \textbf{题量} & \textbf{建议用时} & \textbf{节点控制} \\ \hline
单选题 & 8题 & 25min & 第25分钟开始填空 \\
多选题 & 4题 & 15min & 第40分钟开始解答 \\
填空题 & 4题 & 15min & 第55分钟开始解答17-18 \\
解答题17-18 & 2题 & 20min & 第75分钟进入19-20 \\
解答题19-20 & 2题 & 22min & 第97分钟进入21-22 \\
解答题21-22 & 2题 & 23min & 第120分钟停笔 \\
\end{tabular}
\end{center}

\subsection*{考场铁律}
\begin{enumerate}
  \item \textbf{卡壳3分钟}：一道选择题/填空题思考超过3分钟无思路，标记后跳过
  \item \textbf{多选宁少勿滥}：只选100\%确定的选项，不确定坚决不选（2分vs0分）
  \item \textbf{大题不空题}：写"解" + 关键公式就值2分
  \item \textbf{最后5分钟}：只做两件事——涂答题卡、检查选填是否填涂正确
\end{enumerate}

\newpage

%=================== 试卷正文 ===================
\section*{一、单选题（共8题，每题5分，共40分）}

\textbf{1.} 已知集合 $A=\{x\mid x^2-4x+3<0\}$，$B=\{x\mid 2^x>4\}$，则 $A\cap B=$
\quad A. $(1,2)$ \quad B. $(2,3)$ \quad C. $(1,+\infty)$ \quad D. $(3,+\infty)$

\textbf{2.} 复数 $z=\dfrac{2i}{1+i}$ 的共轭复数为
\quad A. $1-i$ \quad B. $1+i$ \quad C. $-1+i$ \quad D. $-1-i$

\textbf{3.} 已知向量 $\vec{a}=(1,2)$，$\vec{b}=(m,-1)$，若 $\vec{a}\perp\vec{b}$，则 $m=$
\quad A. $2$ \quad B. $-2$ \quad C. $\frac12$ \quad D. $-\frac12$

\textbf{4.} 在 $\triangle ABC$ 中，$a=2$，$b=3$，$\cos C=\frac14$，则 $c=$
\quad A. $\sqrt{10}$ \quad B. $4$ \quad C. $\sqrt{13}$ \quad D. $2\sqrt{3}$

\textbf{5.} 函数 $f(x)=\ln(x^2-2x-3)$ 的单调递增区间是
\quad A. $(1,+\infty)$ \quad B. $(3,+\infty)$ \quad C. $(-\infty,-1)$ \quad D. $(-\infty,1)$

\textbf{6.} 若 $\sin\alpha+\cos\alpha=\frac15$，则 $\sin2\alpha=$
\quad A. $\frac{24}{25}$ \quad B. $-\frac{24}{25}$ \quad C. $\frac{12}{25}$ \quad D. $-\frac{12}{25}$

\textbf{7.} 已知 $F_1,F_2$ 是椭圆 $C:\frac{x^2}{9}+\frac{y^2}{5}=1$ 的两个焦点，$P$ 为 $C$ 上一点，且 $\angle F_1PF_2=60^\circ$，则 $\triangle PF_1F_2$ 的面积为
\quad A. $\frac{5\sqrt3}{3}$ \quad B. $5\sqrt3$ \quad C. $\frac{5\sqrt3}{2}$ \quad D. $5$

\textbf{8.} 设 $a=e^{0.2}$，$b=1.2$，$c=\ln1.2$，则
\quad A. $a<b<c$ \quad B. $c<b<a$ \quad C. $b<c<a$ \quad D. $c<a<b$

\section*{二、多选题（共4题，每题5分，共20分。全部选对得5分，部分选对得2分）}

\textbf{9.} 已知函数 $f(x)=x^3-3x$，则
\quad A. $f(x)$ 是奇函数 \quad B. $f(x)$ 在 $(-\infty,+\infty)$ 上单调递增
\quad C. $f(x)$ 的极大值为 $2$ \quad D. $f(x)$ 在 $[-2,2]$ 上的最小值为 $-2$

\textbf{10.} 下列说法正确的是
\quad A. 命题"$\forall x>0$，$e^x>x+1$"的否定是"$\exists x>0$，$e^x\le x+1$"
\quad B. 若 $a>b$，则 $ac^2>bc^2$
\quad C. "$x>2$"是"$x>1$"的充分不必要条件
\quad D. 函数 $y=2^x$ 与 $y=\log_2 x$ 的图象关于 $y=x$ 对称

\textbf{11.} 等差数列 $\{a_n\}$ 的前 $n$ 项和为 $S_n$，$a_1=1$，$S_3=9$，则
\quad A. $d=2$ \quad B. $a_n=2n-1$ \quad C. $S_n=n^2$ \quad D. $\frac1{a_1a_2}+\frac1{a_2a_3}+\cdots+\frac1{a_8a_9}=\frac89$

\textbf{12.} 已知函数 $f(x)=\begin{cases} |\ln x|, & x>0 \\ 2^{x+1}, & x\le 0 \end{cases}$，若函数 $g(x)=f(x)-k$ 有3个零点，则
\quad A. $k=1$ \quad B. $0<k<1$ \quad C. $k>2$ \quad D. $k=2$

\section*{三、填空题（共4题，每题5分，共20分）}

\textbf{13.} $(x-\frac{2}{x})^5$ 展开式中 $x^3$ 的系数为 \underline{\hspace{2cm}}。

\textbf{14.} 曲线 $y=x^3+2x$ 在点 $(1,3)$ 处的切线方程为 \underline{\hspace{2cm}}。

\textbf{15.} 已知 $\alpha\in(0,\frac{\pi}{2})$，$\cos2\alpha=\frac78$，则 $\tan\alpha=$ \underline{\hspace{2cm}}。

\textbf{16.} 已知 $S_n$ 是数列 $\{a_n\}$ 的前 $n$ 项和，$a_1=2$，$a_{n+1}=2S_n+2$，则 $a_n=$ \underline{\hspace{2cm}}。

\newpage
\section*{四、解答题（共6题，共70分）}

\textbf{17.}（10分）已知数列 $\{a_n\}$ 是等差数列，$a_1=2$，$a_3=6$。
（1）求 $\{a_n\}$ 的通项公式；
（2）设 $b_n=2^{a_n}$，求 $\{b_n\}$ 的前 $n$ 项和 $T_n$。

\vspace{8cm}

\textbf{18.}（12分）在 $\triangle ABC$ 中，角 $A,B,C$ 的对边分别为 $a,b,c$，$2b\cos C=2a-c$。
（1）求 $B$；
（2）若 $b=2\sqrt3$，求 $\triangle ABC$ 面积的最大值。

\vspace{8cm}

\textbf{19.}（12分）如图，四棱锥 $P-ABCD$ 中，底面 $ABCD$ 是矩形，$PA\perp$ 底面 $ABCD$，$PA=AD=2$，$AB=3$，$E$ 为 $PD$ 中点。
（1）求证：$AE\parallel$ 平面 $PBC$；
（2）求二面角 $E-AC-B$ 的余弦值。

\vspace{8cm}

\textbf{20.}（12分）甲、乙两人进行投篮比赛，规则如下：每人投3次，每次命中得2分，未中得0分。甲每次命中概率为 $\frac23$，乙每次命中概率为 $\frac12$。
（1）求甲恰好得4分的概率；
（2）设 $X$ 为甲的全部得分，求 $X$ 的分布列和期望。

\vspace{8cm}

\textbf{21.}（12分）椭圆 $C:\frac{x^2}{a^2}+\frac{y^2}{b^2}=1\;(a>b>0)$ 的离心率为 $\frac12$，右焦点为 $F(1,0)$。
（1）求椭圆 $C$ 的方程；
（2）过点 $P(0,2)$ 的直线 $l$ 与 $C$ 交于 $A,B$，求 $\triangle AOB$ 面积的最大值（$O$ 为原点）。

\vspace{8cm}

\textbf{22.}（12分）已知函数 $f(x)=x\ln x - ax^2 - x$，$a\in\mathbb{R}$。
（1）当 $a=0$ 时，求 $f(x)$ 的极值；
（2）若 $f(x)$ 有两个极值点 $x_1,x_2$，求 $a$ 的取值范围；
（3）求证：当 $a=0$ 时，$f(x) < e^x - x - 1$ 对 $x>0$ 恒成立。

\newpage
%=================== 参考答案 ===================
\section*{参考答案与解析}

\subsection*{一、单选题}
\textbf{1. B} $A=(1,3)$，$B=(2,+\infty)$，$A\cap B=(2,3)$。\\
\textbf{2. B} $z=\frac{2i}{1+i}=\frac{2i(1-i)}{(1+i)(1-i)}=\frac{2+2i}{2}=1+i$，$\bar{z}=1-i$。\\
\textbf{3. A} $\vec{a}\perp\vec{b}\Rightarrow 1\cdot m+2\cdot(-1)=0\Rightarrow m=2$。\\
\textbf{4. A} $c^2=a^2+b^2-2ab\cos C=4+9-12\cdot\frac14=10$，$c=\sqrt{10}$。\\
\textbf{5. B} 定义域 $x^2-2x-3>0\Rightarrow x<-1$ 或 $x>3$。$t=x^2-2x-3$ 在 $(3,+\infty)$ 上增，$y=\ln t$ 增，复合后增。\\
\textbf{6. B} $\sin\alpha+\cos\alpha=\frac15$，平方得 $1+2\sin\alpha\cos\alpha=\frac1{25}$，$\sin2\alpha=-\frac{24}{25}$。\\
\textbf{7. A} $a=3$，$c=2$，$|PF_1|+|PF_2|=6$。余弦定理：
$|F_1F_2|^2=|PF_1|^2+|PF_2|^2-|PF_1||PF_2|$，
$16=(|PF_1|+|PF_2|)^2-3|PF_1||PF_2|=36-3|PF_1||PF_2|$，
$|PF_1||PF_2|=\frac{20}{3}$，$S=\frac12|PF_1||PF_2|\sin60^\circ=\frac12\cdot\frac{20}{3}\cdot\frac{\sqrt3}{2}=\frac{5\sqrt3}{3}$。\\
\textbf{8. B} $e^{0.2}>1.2$（$e^{0.2}\approx1.221$），$\ln1.2<0.2$（$\ln1.2\approx0.182$），$a>1.2>b>0.2>c$。

\subsection*{二、多选题}
\textbf{9. AC} $f(-x)=-f(x)$ 奇函数；$f'(x)=3x^2-3=3(x-1)(x+1)$，$f_{\max}=f(-1)=2$。\\
\textbf{10. ACD} B错误：$c=0$时$ac^2=bc^2$。\\
\textbf{11. ABC} $S_3=3a_1+3d=3+3d=9\Rightarrow d=2$，$a_n=2n-1$，$S_n=n^2$，$\sum_{i=1}^{8}\frac1{(2i-1)(2i+1)}=\frac12(1-\frac1{17})=\frac8{17}$。\\
\textbf{12. BD} 画图知$0<k<1$有3个零点，$k=2$有3个零点（$x=-1$和两个正根）。

\subsection*{三、填空题}
\textbf{13.} $T_{k+1}=C_5^k x^{5-k}(-\frac{2}{x})^k=C_5^k(-2)^k x^{5-2k}$，令$5-2k=3\Rightarrow k=1$，系数为$C_5^1(-2)^1=-10$。\\
\textbf{14.} $y'=3x^2+2$，$k=y'(1)=5$，$y-3=5(x-1)\Rightarrow y=5x-2$。\\
\textbf{15.} $\cos2\alpha=1-2\sin^2\alpha=\frac78\Rightarrow\sin^2\alpha=\frac1{16}$，$\sin\alpha=\frac14$，$\cos\alpha=\frac{\sqrt{15}}{4}$，$\tan\alpha=\frac1{\sqrt{15}}$。 \\
\textbf{16.} $a_{n+1}=2S_n+2$，$a_n=2S_{n-1}+2$，相减得$a_{n+1}=3a_n$，$a_2=2S_1+2=6$，$n\ge2$时$a_n=6\cdot3^{n-2}=2\cdot3^{n-1}$。

\subsection*{四、解答题}
\textbf{17.}（1）$d=\frac{a_3-a_1}{2}=2$，$a_n=2+(n-1)\cdot2=2n$。\\
（2）$b_n=2^{2n}=4^n$，$T_n=\frac{4(4^n-1)}{4-1}=\frac43(4^n-1)$。

\textbf{18.}（1）略（同前，$B=60^\circ$）。\\
（2）$S=\frac12ac\sin B$，$b^2=a^2+c^2-2ac\cos B\Rightarrow12=a^2+c^2-ac\ge2ac-ac=ac$，$ac\le12$，$S_{\max}=\frac12\cdot12\cdot\frac{\sqrt3}{2}=3\sqrt3$。

\textbf{19.}（1）取$PB$中点$F$，证$EF\parallel AD\parallel BC$，$AEFD$平行四边形$\Rightarrow AE\parallel DF\subset$面$PBC$。\\
（2）建系$A(0,0,0),B(3,0,0),D(0,2,0),P(0,0,2),E(0,1,1)$，$\vec{n_1}=(0,0,1)$（面$ABC$法向量），$\vec{n_2}=(2,3,0)$（面$AEC$法向量），$\cos\theta=\frac{|\vec{n_1}\cdot\vec{n_2}|}{|\vec{n_1}||\vec{n_2}|}=\frac{3}{\sqrt{13}}$。

\textbf{20.}（1）甲投3次命2次：$P=C_3^2(\frac23)^2(\frac13)=\frac49$。\\
（2）$X\sim B(3,\frac23)$，$P(X=k)=C_3^k(\frac23)^k(\frac13)^{3-k}$，
$E(X)=3\cdot\frac23=2$。

\textbf{21.}（1）$e=\frac12$，$c=1$，$a=2$，$b^2=a^2-c^2=3$，$\frac{x^2}{4}+\frac{y^2}{3}=1$。\\
（2）设$l:y=kx+2$，联立得$(3+4k^2)x^2+16kx+4=0$，
$\Delta=256k^2-16(3+4k^2)=192k^2-48>0\Rightarrow k^2>\frac14$。
$S_{\triangle AOB}=\frac12\cdot2\cdot|x_1-x_2|=|x_1-x_2|=\frac{\sqrt{192k^2-48}}{3+4k^2}=\frac{4\sqrt3\sqrt{4k^2-1}}{4k^2+3}$。
令$t=\sqrt{4k^2-1}>0$，$S=4\sqrt3\cdot\frac{t}{t^2+4}\le4\sqrt3\cdot\frac1{2\sqrt4}=4\sqrt3\cdot\frac14=\sqrt3$（$t=2$时取等）。

\textbf{22.}（1）$a=0$，$f(x)=x\ln x-x$，$f'(x)=\ln x$，
$(0,1)$减$(1,+\infty)$增，极小值$f(1)=-1$，无极大值。\\
（2）$f'(x)=\ln x-2ax$有两个变号零点$\Rightarrow$...（完整推导略）。\\
（3）令$g(x)=e^x-x-1-x\ln x+x=e^x-x\ln x-1$，
$g'(x)=e^x-\ln x-1$，$g''(x)=e^x-\frac1x>0$，$g'(1)=e-1>0$，
$g'(x)$增且$g'(1)>0$，$g(x)$增，$g(0^+)\to0$，$\therefore g(x)>0$。

\vfill
\centerline{\textbf{—— 本卷完 ——}}
\end{document}
